\section{Introduction}
\label{sec:introduction}

Efficient resource allocation is critical in modern manufacturing, transportation, and distribution systems \citep{pinedo2016}.
While classical shop scheduling models focus on minimizing total completion time (makespan) or tardiness, modern production systems increasingly operate under Just-In-Time (JIT) principles.
Under JIT scheduling, completing a task too early increases inventory holding costs, whereas late completion incurs tardiness penalties and compromises customer satisfaction.
Consequently, JIT scheduling requires minimizing both earliness and tardiness penalties simultaneously.

The Just-In-Time Job Shop Scheduling (JITJSS) problem \citep{baptiste2008lagrangian} formally models this operational context.
In JITJSS, each job consists of operations with specific target due dates, and the objective is to determine an optimal machine schedule that minimizes total weighted earliness and tardiness costs.
JITJSS is a strongly NP-hard problem \citep{du1990minimizing}.
Consequently, extensive research has focused on metaheuristic algorithms, including genetic algorithms \citep{araujo2009genetic, dos2010combination, yang2012improved}, variable neighborhood search \citep{wang2014variable, ahmadian2021meta}, adaptive large neighborhood search \citep{laborie2007self, abderrazzak2022adaptive}, and deep reinforcement learning \citep{sabri2023reinforcement}.
Across the JITJSS literature, five key studies \citep{dos2010combination, wang2014variable, ahmadian2021meta, abderrazzak2022adaptive, sabri2023reinforcement} hold the best-known solution value for at least one benchmark instance.
Among them, \citeauthor{ahmadian2021meta} provides the most comprehensive baseline by evaluating all problem scales, whereas other recent works \citep{abderrazzak2022adaptive, sabri2023reinforcement} focused on small and medium instances.
Evaluating algorithmic performance across the benchmark set therefore requires comparing solutions against the composite best results established across these five studies.

In JITJSS, empirical evidence reveals that tardiness penalties dominate total schedule costs.
However, existing neighborhood operators in the literature do not focus on this critical aspect.
Most published methods rely either on random, unfiltered operation moves \citep{dos2010combination, ahmadian2021meta}, or on local search operators adapted from makespan minimization ignoring operation due dates.
Consequently, standard operators lack mechanisms to identify and resolve the specific operation sequences that delay tardy jobs.
To address this gap, this paper introduces a novel neighborhood operator, named \texttt{swap\_crit}, specifically designed to mitigate tardiness penalties.
The \texttt{swap\_crit} operator traces the backward critical path of each individual tardy operation, identifies the exact sequence of bottleneck operations causing deadline violations, and generates candidate moves by swapping boundary operations of critical blocks along that path.

A key methodological challenge in metaheuristic research is evaluating the true merit of a newly proposed operator without designer bias.
Traditionally, researchers embed a new operator into a fixed metaheuristic framework and tune its parameters manually \citep{birattari2009tuning}.
This manual approach risks overfitting the operator to a specific algorithmic architecture and fails to demonstrate whether the operator would remain effective in alternative configurations \citep{stutzle2018automated}.
To ensure an unbiased and rigorous validation, we employ Automatic Algorithm Design (AAD) via \texttt{irace} \citep{lopez2016irace} as an experimental testing tool.
We formalize a broad algorithmic design space using a context-free grammar in Backus-Naur Form (BNF) \citep{mascia2013grammars}.
This design space couples our proposed \texttt{swap\_crit} operator with alternative literature neighborhoods (classical swaps, penalty heuristics, and exact mathematical relaxations), diverse initialization procedures, multiple search frameworks (Local Search, Tabu Search, and Iterated Local Search), and sequence schedulers.
By allowing the automated configurator to select components freely, we directly test whether \texttt{swap\_crit} is selected over competing literature operators and whether it contributes decisively to achieving state-of-the-art performance.

The main contributions of this paper are fourfold:
\begin{enumerate}[(i)]
    \item We propose \texttt{swap\_crit}, a novel neighborhood operator for JITJSS that traces individual backward critical paths to resolve bottleneck operations causing tardiness penalties.
    \item We formulate a metaheuristic search that integrates \texttt{swap\_crit} within a multi-neighborhood Iterated Tabu Search (ITS) architecture.
    \item We establish an AAD validation framework parameterized via a context-free BNF grammar to evaluate the proposed operator objectively within diverse algorithm components.
    \item We conduct extensive computational experiments on 72 standard benchmark instances \citep{baptiste2008lagrangian}, showing that \texttt{irace} selects \texttt{swap\_crit} in 80\% of elite configurations, and the resulting algorithm matches or improves upon literature results in 72.2\% of the instances.
\end{enumerate}

The remainder of this paper is organized as follows:
Section~\ref{sec:swap_crit} introduces the proposed \texttt{swap\_crit} neighborhood operator.
Section~\ref{sec:design_space} defines the algorithmic design space and its context-free grammar.
Section~\ref{sec:config_process} describes the automated algorithm configuration protocol.
Section~\ref{sec:results} presents computational experiments and comparative analyses against literature baselines.
Section~\ref{sec:conclu} concludes the paper with final remarks.
